\chapter{핵심 수식 유도}
\label{app:derivations}

\section{softmax의 shift invariance와 안정적 계산}

softmax는 모든 logit에 같은 상수를 더해도 변하지 않는다. 임의의 $c\in\R$에 대해
\begin{equation}
  \frac{e^{z_i+c}}{\sum_j e^{z_j+c}}
  =\frac{e^c e^{z_i}}{e^c\sum_j e^{z_j}}
  =\frac{e^{z_i}}{\sum_j e^{z_j}}
\end{equation}
가 성립하기 때문이다.

안정적인 softmax는 가장 큰 logit을 빼서 overflow를 줄인다. $m=\max_j z_j$로 두면
\begin{equation}
  p_i=\frac{e^{z_i-m}}{\sum_j e^{z_j-m}}
\end{equation}
이고 모든 지수의 입력이 0 이하가 된다. 이 변환은 위의 shift invariance 때문에 확률을
바꾸지 않는다.

softmax의 Jacobian은 vocabulary logit 사이의 결합을 보여 준다. 미분하면
\begin{equation}
  \frac{\partial p_i}{\partial z_j}=p_i(\mathbf{1}[i=j]-p_j)
\end{equation}
를 얻는다. 따라서 token $j$의 logit 변화는 $j$의 확률뿐 아니라 다른 모든 token의
확률에도 영향을 준다.

\section{cross-entropy와 logit gradient}

one-hot target $y$에 대한 softmax cross-entropy의 logit gradient는 $\vect{p}-\vect{y}$다.
target index를 $k$라고 하면
\begin{align}
  \mathcal{L}&=-\log p_k
  =-z_k+\log\sum_j e^{z_j},\\
  \frac{\partial\mathcal{L}}{\partial z_i}
  &=-\mathbf{1}[i=k]+\frac{e^{z_i}}{\sum_j e^{z_j}}
  =p_i-y_i.
\end{align}
이 식은 정답 logit을 올리고 현재 확률에 비례하여 다른 logit을 내리는 gradient 신호를
보여 준다.

LM head의 weight gradient는 hidden state와 logit error의 outer product다. row-vector
convention에서 $\vect{z}=\mat{W}_U\vect{h}+\vect{b}$라면
\begin{equation}
  \frac{\partial\mathcal{L}}{\partial\mat{W}_U}
  =(\vect{p}-\vect{y})\vect{h}^{\mathsf T},\qquad
  \frac{\partial\mathcal{L}}{\partial\vect{h}}
  =\mat{W}_U^{\mathsf T}(\vect{p}-\vect{y}).
\end{equation}
이 gradient는 chain rule을 통해 앞 block과 embedding까지 전달된다.

\section{scaled dot-product attention의 차원과 scale}

attention score matrix의 두 $T$축은 destination과 source position에서 생긴다. batch와
head 축을 생략하면
\begin{equation}
  \mat{Q}\in\R^{T\times d_h},\quad
  \mat{K}^{\mathsf T}\in\R^{d_h\times T},\quad
  \mat{Q}\mat{K}^{\mathsf T}\in\R^{T\times T}
\end{equation}
이며, $\mat{A}\mat{V}$는
$(T\times T)(T\times d_h)=T\times d_h$가 된다.

$1/\sqrt{d_h}$ scale은 독립이고 평균 0, 분산 1인 query/key coordinate라는 이상화에서
dot-product 분산을 차원과 무관하게 만든다. $q_r k_r$들이 독립이고 분산 1이면
\begin{equation}
  \Var\!\left(\sum_{r=1}^{d_h}q_r k_r\right)=d_h,\qquad
  \Var\!\left(\frac{1}{\sqrt{d_h}}\sum_{r=1}^{d_h}q_r k_r\right)=1.
\end{equation}
실제 학습된 q와 k는 이 독립 가정을 정확히 만족하지 않지만, 유도는 scale의 설계 동기를
설명한다 \citep{vaswani2017attention}.

causal mask는 softmax 뒤가 아니라 score에 적용하는 것이 정확한 정의다. 미래 score를
$-\infty$로 두면 그 지수는 0이 되고 허용된 source만 다시 정규화된다. softmax 뒤의
weight를 0으로 만들고 재정규화하지 않으면 row sum이 1이 아니므로 다른 함수가 된다.

\section{여러 head의 출력은 residual update의 합이다}

concatenation 뒤의 output projection은 head별 output projection의 합으로 분해된다.
$\mat{O}_{\mathrm{cat}}=[\mat{O}^{1}\;\cdots\;\mat{O}^{H}]$이고
$\mat{W}_O$를 대응하는 row block으로
\begin{equation}
  \mat{W}_O=
  \begin{bmatrix}
    \mat{W}_O^1\\ \vdots\\ \mat{W}_O^H
  \end{bmatrix}
\end{equation}
라고 나누면 block multiplication으로
\begin{equation}
  \mat{O}_{\mathrm{cat}}\mat{W}_O
  =\sum_{h=1}^{H}\mat{O}^{h}\mat{W}_O^h
\end{equation}
를 얻는다. 이 등식은 head output을 residual 방향별로 분석하는 근거다.

head별 합은 head가 기능적으로 독립이라는 뜻이 아니다. 모든 head의 Q/K/V가 같은 residual
input을 읽고, 뒤 layer가 합쳐진 결과를 비선형적으로 처리하므로 한 head의 효과는 다른
head와 경로에 의존할 수 있다.

\section{RoPE가 상대 offset을 포함하는 이유}

2차원 회전행렬은 $\mat{R}(a)^{\mathsf T}\mat{R}(b)=\mat{R}(b-a)$를 만족한다. 삼각함수의
덧셈정리를 회전행렬 곱에 적용하면 바로 확인할 수 있다. 따라서
\begin{equation}
  (\mat{R}(i)\vect{q})^{\mathsf T}(\mat{R}(j)\vect{k})
  =\vect{q}^{\mathsf T}\mat{R}(j-i)\vect{k}
\end{equation}
이고 score가 절대 position의 차이 $j-i$에 의존한다.

RoPE의 상대성은 content vector와 offset의 상호작용이다. 같은 offset이라도 $\vect{q}$와
$\vect{k}$가 다르면 score가 다르며, RoPE만으로 특정 distance를 선택하는 attention pattern이
보장되지는 않는다. frequency schedule과 학습된 projection이 함께 작동한다.

\section{RMSNorm이 보존하고 바꾸는 양}

gain과 $\epsilon$을 잠시 제외한 RMS normalization은 출력의 root-mean-square를 1로 만든다.
$r(\vect{x})=\sqrt{d^{-1}\sum_i x_i^2}$라면
\begin{equation}
  r\!\left(\frac{\vect{x}}{r(\vect{x})}\right)=1
\end{equation}
이다. 실제 식에서는 $\epsilon>0$과 coordinate-wise gain $\vect{g}$ 때문에 최종 output의
RMS가 정확히 1일 필요는 없다.

RMSNorm은 입력의 전체 positive scale에 거의 불변이다. $\epsilon=0$이고 $a>0$이면
\begin{equation}
  \operatorname{RMSNorm}(a\vect{x})
  =\vect{g}\odot\frac{a\vect{x}}{a r(\vect{x})}
  =\operatorname{RMSNorm}(\vect{x}).
\end{equation}
$a<0$이면 방향 부호가 바뀌며, 유한 $\epsilon$에서는 정확한 scale invariance가 아니다.

pre-norm residual stream의 norm은 layer를 거치며 계속 변할 수 있다. RMSNorm은 sublayer가
읽는 사본을 정규화할 뿐, identity path의 residual state를 정규화된 값으로 교체하지 않기
때문이다.

\section{KV cache가 같은 attention을 계산하는 이유}

KV cache는 과거 position의 key와 value가 새 token 때문에 바뀌지 않는 causal decoder에서
projection을 재사용한다. 새 position $T+1$의 hidden state는 과거 position의 hidden state를
바꾸지 않으므로, 같은 layer의 $\mat{K}_{1:T}$와 $\mat{V}_{1:T}$도 이전 prefill/decode 때
계산한 값과 같다.

cache와 전체 재계산의 새-position attention row는 이상적인 산술에서 동일하다. 두 방법
모두 같은 $\vect{q}_{T+1}$과 연결된
$[\mat{K}_{1:T};\vect{k}_{T+1}]$, $[\mat{V}_{1:T};\vect{v}_{T+1}]$를 사용하기 때문이다.
과거 query row는 다음-token logit에 필요하지 않으므로 다시 만들지 않는다.

KV cache의 원소 수는 batch, layer, KV head, sequence, head dimension과 K/V 두 종류의 곱이다.
padding과 metadata를 제외한 단순 원소 수는
\begin{equation}
  N_{\mathrm{cache}}=2BLH_{kv}Td_h
\end{equation}
이며, element당 byte 수를 곱하면 기본 저장량을 얻는다. tensor parallel 복제, alignment,
quantization과 paging은 실제 메모리를 바꾼다.

\section{GQA가 줄이는 parameter와 cache}

GQA는 query projection보다 key/value projection의 output dimension을 줄인다. bias를
제외하고 input dimension이 $d$이면
\begin{align}
  N_Q&=dH_qd_h,\\
  N_K+N_V&=2dH_{kv}d_h
\end{align}
이다. MHA의 $H_{kv}=H_q$에서 GQA의 작은 $H_{kv}$로 바꾸면 K/V parameter와 cache가 같은
비율로 줄지만 Q와 output projection은 이 식에서 줄지 않는다.

GQA의 계산량 감소는 단계에 따라 다르다. K/V projection과 cache traffic은 줄지만 각
query head는 여전히 attention weight와 value sum을 계산한다. kernel과 sequence 길이를
고려하지 않고 전체 inference FLOPs가 $H_{kv}/H_q$ 배로 줄었다고 말하면 부정확하다.

\section{direct logit attribution과 final norm}

final normalization이 없으면 residual update의 logit 기여는 정확히 선형이다. target
token $a$와 contrast token $b$의 unembedding 차이를
$\vect{w}_{ab}=\vect{w}_{U,a}-\vect{w}_{U,b}$라고 하면
\begin{equation}
  z_a-z_b=\vect{w}_{ab}^{\mathsf T}
  \left(\vect{x}^{(0)}+\sum_c\Delta\vect{x}_c\right)+(b_a-b_b)
\end{equation}
이므로 각 $c$의 항은 $\vect{w}_{ab}^{\mathsf T}\Delta\vect{x}_c$다.

final normalization이 있으면 구성요소 기여는 일반적으로 서로 결합된다. 예를 들어
RMSNorm의 denominator는 모든 residual component의 합으로 결정되므로 한 update를 바꾸면
다른 component가 받는 공통 scale도 바뀐다. 연구자는 normalization을 실제로 적용한
component-wise projection, frozen-scale approximation 또는 local Jacobian 가운데 무엇을
사용했는지 밝혀야 한다.

\section{activation patching의 first-order 근사}

attribution patching은 finite intervention을 Taylor 전개의 첫 항으로 근사한다. corrupted
activation을 $\vect{h}_r$, clean activation을 $\vect{h}_c$라고 하면
\begin{equation}
  m(\vect{h}_c)-m(\vect{h}_r)
  =\nabla m(\vect{h}_r)^{\mathsf T}(\vect{h}_c-\vect{h}_r)
  +\mathcal{O}(\lVert\vect{h}_c-\vect{h}_r\rVert^2)
\end{equation}
다. 차이가 작고 곡률이 약할수록 first-order score가 실제 patch effect에 가까울 수 있다.

큰 patch와 activation threshold는 선형 근사의 오차를 키울 수 있다. ReLU나 gate의 on/off가
바뀌거나 attention softmax가 크게 이동하면 고차항이 중요해진다. 계산량을 줄이기 위해
근사를 사용한 뒤 상위 후보에는 실제 patch를 수행하는 이유가 여기에 있다.

\section{Jensen--Shannon divergence}

Jensen--Shannon(JS) divergence는 두 확률분포의 대칭적인 차이를 측정한다. 분포
$P,Q$와 $M=(P+Q)/2$에 대해
\begin{equation}
  \operatorname{JS}(P,Q)=\frac12\KL(P\|M)+\frac12\KL(Q\|M)
\end{equation}
로 정의한다. 자연로그를 쓰면 값은 $0$과 $\log 2$ 사이에 있으며, 로그 base 2를 쓰면
상한은 1이다.

JS divergence는 어떤 token의 logit이 변했는지를 혼자 알려 주지 않는다. vocabulary 전체
분포의 변화를 요약하므로 target token과 무관한 mass 이동도 값에 들어간다. 사실 회상
개입에서는 target logit difference와 JS divergence를 함께 보면 질문이 다른 두 metric을
구분할 수 있다.

\begin{studybox}
각 유도는 작은 무작위 tensor를 double precision으로 계산하여 등식을 검산한다. 수학적
등식의 허용 오차와 GPU mixed-precision kernel의 허용 오차를 따로 정하며, shape와 축을
assertion으로 고정한다.
\end{studybox}
